\documentclass[12pt, a4paper]{article}

\usepackage{xurl}
\usepackage[utf8]{inputenc}
\usepackage[spanish]{babel}
\usepackage{geometry}
\usepackage{titlesec}
\usepackage{enumitem}
\usepackage{xcolor}
\usepackage{listings}
\usepackage{mdframed}
\usepackage{hyperref}
\usepackage{tikz}
\usepackage{graphicx}
\usepackage{float}
\usetikzlibrary{shapes.geometric, arrows.meta, positioning}
 
\geometry{margin=2.5cm}
 
% Colores
\definecolor{azulprincipal}{RGB}{30, 90, 160}
\definecolor{amarillopendiente}{RGB}{255, 200, 0}
\definecolor{verdesprintold}{RGB}{0, 140, 60}
\definecolor{azulsprintnew}{RGB}{0, 80, 180}
\definecolor{rojocambio}{RGB}{180, 30, 30}
\definecolor{grisclaro}{RGB}{245, 245, 245}
\definecolor{codigofondo}{RGB}{240, 244, 250}
\definecolor{codigoborde}{RGB}{30, 90, 160}
\definecolor{fondoverde}{RGB}{230, 255, 235}
\definecolor{fondobordoverde}{RGB}{0, 140, 60}
\definecolor{fondoazul}{RGB}{230, 240, 255}
\definecolor{fondobordoazul}{RGB}{0, 80, 180}
\definecolor{fondoamarillo}{RGB}{255, 252, 220}
\definecolor{fondobordoamarillo}{RGB}{180, 140, 0}
 
% Estilo de secciones
\titleformat{\section}{\large\bfseries\color{azulprincipal}}{}{0em}{}[\titlerule]
\titleformat{\subsection}{\normalsize\bfseries\color{azulprincipal}}{}{0em}{}
 
% Estilo de código
\lstset{
    basicstyle=\ttfamily\small,
    backgroundcolor=\color{codigofondo},
    frame=single,
    rulecolor=\color{codigoborde},
    breaklines=true,
    language=Java,
    keywordstyle=\color{azulprincipal}\bfseries,
    commentstyle=\color{gray},
    stringstyle=\color{teal},
}
 
% Caja Sprint 1 (verde)
\newmdenv[
    backgroundcolor=fondoverde,
    linecolor=fondobordoverde,
    leftmargin=0, rightmargin=0,
    innerleftmargin=10pt, innerrightmargin=10pt,
    innertopmargin=8pt, innerbottommargin=8pt,
    roundcorner=4pt,
    frametitle={\color{verdesprintold}\bfseries Sprint 1 --- Versi\'{o}n inicial},
    frametitleaboveskip=4pt,
    frametitlebelowskip=4pt
]{cajasprint1}
 
% Caja Sprint 2 (amarillo)
\newmdenv[
    backgroundcolor=fondoamarillo,
    linecolor=fondobordoamarillo,
    leftmargin=0, rightmargin=0,
    innerleftmargin=10pt, innerrightmargin=10pt,
    innertopmargin=8pt, innerbottommargin=8pt,
    roundcorner=4pt,
    frametitle={\color{fondobordoamarillo}\bfseries Sprint 2 --- Actualizaci\'{o}n},
    frametitleaboveskip=4pt,
    frametitlebelowskip=4pt
]{cajasprint2}
 
% Caja versión final (azul)
\newmdenv[
    backgroundcolor=fondoazul,
    linecolor=fondobordoazul,
    leftmargin=0, rightmargin=0,
    innerleftmargin=10pt, innerrightmargin=10pt,
    innertopmargin=8pt, innerbottommargin=8pt,
    roundcorner=4pt,
    frametitle={\color{azulsprintnew}\bfseries Versi\'{o}n Final},
    frametitleaboveskip=4pt,
    frametitlebelowskip=4pt
]{cajafinal}
 
% --------------------------------------------------
\begin{document}
 
\begin{center}
    {\LARGE \textbf{Documentaci\'{o}n: Arquitectura y Patrones de Dise\~{n}o}}\\[0.4em]
    {\large \textit{Space Invaders --- Proyecto de Programaci\'{o}n}}\\[0.3em]
    {\small Fecha: \today \quad | \quad Grupo: Cachopin}
\end{center}
 
\vspace{0.5em}
\hrule
\vspace{1.5em}
 
\tableofcontents
\newpage
 
% ==================================================
\section{Introducci\'{o}n General}
 
En este proyecto se ha implementado una arquitectura basada en dos enfoques complementarios: el patr\'{o}n arquitect\'{o}nico \textbf{MVC (Modelo-Vista-Controlador)} combinado con el patr\'{o}n \textbf{Observer} para la sincronizaci\'{o}n entre componentes, y la aplicaci\'{o}n de tres patrones de dise\~{n}o (\textbf{Factory}, \textbf{Strategy} y \textbf{Composite}) para resolver problemas espec\'{i}ficos de flexibilidad y mantenibilidad.
 
El proyecto ha pasado por una evoluci\'{o}n clara entre el \textbf{Sprint 1}, donde se implement\'{o} una versi\'{o}n funcional pero b\'{a}sica, el \textbf{Sprint 2}, donde se refactoriz\'{o} profundamente el modelo para soportar m\'{u}ltiples naves, disparos compuestos, flota de enemigos y detecci\'{o}n de colisiones a nivel de p\'{i}xel, y la \textbf{versi\'{o}n final}, la versi\'{o}n actual y definitiva del juego. Esta documentaci\'{o}n refleja todas las etapas y el razonamiento detr\'{a}s de cada cambio.
 
\newpage
 
% ==================================================
\part*{Parte I: Arquitectura MVC + Observer}
\addcontentsline{toc}{section}{Parte I: Arquitectura MVC + Observer}
 
% ==================================================
\section{Introducci\'{o}n al Patr\'{o}n MVC + Observer}
 
El patr\'{o}n \textbf{Modelo-Vista-Controlador (MVC)} es una arquitectura de dise\~{n}o utilizada en programaci\'{o}n orientada a objetos que permite organizar una aplicaci\'{o}n separando sus componentes en tres partes principales. Esta separaci\'{o}n facilita el mantenimiento del c\'{o}digo, mejora su organizaci\'{o}n y permite que cada componente tenga una responsabilidad clara dentro del sistema.
 
\begin{description}[leftmargin=2em, style=nextline]
    \item[\textbf{Modelo}]
        Se encarga de gestionar los datos y la l\'{o}gica del programa. Contiene toda la l\'{o}gica de negocio del juego y las operaciones necesarias para modificarla o consultarla. El modelo no depende de la interfaz gr\'{a}fica ni de c\'{o}mo se muestran los datos al usuario.
 
    \item[\textbf{Vista}]
        Es la parte encargada de presentar la informaci\'{o}n al usuario. Su funci\'{o}n principal es mostrar los datos que proporciona el modelo y actualizarse cuando estos cambian. La vista no contiene la l\'{o}gica principal del programa, sino \'{u}nicamente la representaci\'{o}n visual de la informaci\'{o}n.
 
    \item[\textbf{Controlador}]
        Act\'{u}a como intermediario entre el usuario y el sistema. Recibe las acciones del usuario (clics, entradas de teclado, etc.), interpreta estas acciones y solicita al modelo que realice las operaciones correspondientes.
\end{description}
 
A esta arquitectura se le a\~{n}ade el patr\'{o}n \textbf{Observer (Observador)} para que la vista se mantenga actualizada cuando cambian los datos del modelo. Este patr\'{o}n establece una relaci\'{o}n entre un objeto denominado \textit{observable} (el modelo) y uno o varios \textit{observadores} (las vistas). Cuando el estado del observable cambia, todos los observadores registrados son notificados autom\'{a}ticamente y pueden actualizar su informaci\'{o}n.
 
La combinaci\'{o}n de MVC y Observer permite desarrollar aplicaciones m\'{a}s modulares, donde la l\'{o}gica del programa, la presentaci\'{o}n de los datos y la interacci\'{o}n con el usuario est\'{a}n claramente separadas. Esto facilita la reutilizaci\'{o}n del c\'{o}digo, la ampliaci\'{o}n del sistema y el mantenimiento a largo plazo.
 
% ==================================================
\section{Implementaci\'{o}n en Space Invaders}
 
\subsection{La Capa Modelo}
 
La capa modelo es la que mayor evoluci\'{o}n ha sufrido entre sprints. A continuaci\'{o}n se describe c\'{o}mo era en cada etapa y qu\'{e} motiv\'{o} cada cambio.
 
% ---- SPRINT 1 ----
\begin{cajasprint1}
La clase central era \texttt{Espacio} (Singleton y Observable). Dentro de ella viv\'{i}a toda la l\'{o}gica: el \texttt{game loop} con un \texttt{Timer} de 50ms, la detecci\'{o}n de colisiones, y las condiciones de victoria y derrota.
 
Las clases auxiliares eran simples y directas:
\begin{itemize}
    \item \texttt{Naves} --- clase abstracta con \texttt{x}, \texttt{y}, \texttt{velocidad} y \texttt{vivo}. Defin\'{i}a el m\'{e}todo abstracto \texttt{mover(int dx, int dy)}.
    \item \texttt{Jugador} --- extend\'{i}a \texttt{Naves}. Ten\'{i}a un \'{u}nico objeto \texttt{Disparo} asociado. Al disparar, creaba un nuevo \texttt{Disparo} en la posici\'{o}n \texttt{(x, y-1)}. Solo un disparo activo a la vez.
    \item \texttt{Enemigo} --- extend\'{i}a \texttt{Naves}. Implementaba \texttt{mover()} sumando directamente el desplazamiento. Un \'{u}nico enemigo inicial en posici\'{o}n aleatoria de la fila superior.
    \item \texttt{Disparo} --- clase simple con \texttt{x}, \texttt{y} y \texttt{activo}. El m\'{e}todo \texttt{subir()} decrementaba Y un p\'{i}xel por tick. Se desactivaba al salir del tablero.
\end{itemize}
 
La detecci\'{o}n de colisiones en \texttt{Espacio} comparaba la posici\'{o}n del disparo con la del enemigo con una ventana de 2 p\'{i}xeles para compensar la discretizaci\'{o}n:
\begin{lstlisting}
if (e.isVivo() && d.getX() == e.getX()
        && d.getY() <= e.getY() && d.getY() >= e.getY() - 1) {
    e.setVivo(false);
    d.setActivo(false);
    ...
}
\end{lstlisting}
 
Todo era sencillo y funcionaba, pero el modelo estaba muy acoplado: \texttt{Espacio} sab\'{i}a demasiado sobre \texttt{Jugador}, \texttt{Enemigo} y \texttt{Disparo}, y cualquier extensi\'{o}n requer\'{i}a modificar directamente \texttt{Espacio}.
\end{cajasprint1}
 
\vspace{0.8em}
\textbf{¿Qu\'{e} limitaciones motivaron el cambio al Sprint 2?} Al finalizar el primer Sprint y despu\'{e}s de la reuni\'{o}n con el profesor, buscamos arreglar los fallos comentados, como que solo apareciese un enemigo. Adem\'{a}s, al querer a\~{n}adir m\'{u}ltiples tipos de nave, disparos con formas distintas, varios enemigos simult\'{a}neos con movimiento independiente y colisiones precisas p\'{i}xel a p\'{i}xel, el modelo del Sprint 1 no escalaba. Cada nueva funcionalidad obligaba a tocar \texttt{Espacio}, rompiendo el principio de responsabilidad \'{u}nica.
 
\vspace{0.5em}
 
% ---- SPRINT 2 ----
\begin{cajasprint2}
En este nuevo Sprint se produjo un cambio profundo en la arquitectura. El controlador transmitir\'{i}a los cambios del juego a una nueva clase \texttt{JugadorBueno}, y se introdujeron nuevos patrones que modificaron el dise\~{n}o de las clases:
 
\begin{itemize}
    \item \texttt{Espacio} (Singleton + Observable) --- Evolucion\'{o} su papel en el proyecto: pas\'{o} de ser la clase que gest\'{i}onaba toda la l\'{o}gica a ser la que notificar\'{i}a a la vista de los cambios l\'{o}gicos del juego. En este Sprint se mantuvieron algunos aspectos clave como el Timer de 50ms y las condiciones de victoria/derrota, y segu\'{i}a siendo la encargada de inicializar la nave y los enemigos mediante llamadas directas a \texttt{JugadorBueno} y \texttt{FlotaEnemigos}.
    \item \texttt{JugadorBueno} (Singleton + Observer) --- Nueva clase encargada de recibir la informaci\'{o}n del controlador y enviar se\~{n}ales a las diferentes clases para cambiar la l\'{o}gica del juego. Se responsabilizaba de la construcci\'{o}n de la nave inicial y los cambios de disparo, aunque todav\'{i}a no enviaba se\~{n}ales de inicializaci\'{o}n propias.
    \item \texttt{FlotaEnemigos} (Singleton) --- Nueva clase encargada de almacenar a todos los enemigos del juego, gestionando su movimiento y las colisiones con los disparos del jugador.
    \item \textbf{\texttt{Naves}} --- Sigue siendo abstracta, pero ahora gestiona un \texttt{Composite} (\'{a}rbol de p\'{i}xeles que forma la nave) y un \texttt{Disparo} (gestor de estrategias). Define el m\'{e}todo abstracto \texttt{construir()} con el template method para que cada subclase dibuje su propia geometr\'{i}a.
    \item \textbf{\texttt{Nave1}, \texttt{Nave2}, \texttt{Nave3}} --- Subclases concretas de \texttt{Naves}. Cada una define su geometr\'{i}a de p\'{i}xeles en \texttt{construir()} y sus estrategias de disparo disponibles.
    \item \textbf{\texttt{Enemigo}} --- Ahora es multip\'{i}xel (tambi\'{e}n usa el \texttt{Composite}).
    
\end{itemize}
\end{cajasprint2}
 
\vspace{0.8em}
\textbf{¿Qu\'{e} limitaciones motivaron el cambio al Sprint 2?} 
Al finalizar el segundo Sprint, enocntramos el siguiente problema: Espacio no podia comunicarse directamente con otras clases. Sin embargo, nos interesaba que se comunicase con FlotaEnemigos y JuagdorBueno, por lo que decidimos implementar el patron Observer y, de esa forma, puedan comunicarse. Además, esto significaba que la inicialización se haria mediante notificaciones, y que desde Espacio necesitariamos conocer en todo momento las posiciones de todas las naves y disparos, de ahi surgio la idea de crear una matriz que guardase dichas posiciones.
 
\vspace{0.5em}

% ---- VERSIÓN FINAL ----
\begin{cajafinal}
La versi\'{o}n final redistribuye las responsabilidades de forma mucho m\'{a}s granular:
 
\begin{itemize}
    \item \textbf{\texttt{Espacio}} (Singleton + Observable) --- Sigue siendo el n\'{u}cleo observable, pero ya no contiene el game loop ni la l\'{o}gica de movimiento. Su funci\'{o}n principal es mantener una \textbf{matriz espejo} (\texttt{int[100][60]}) que registra qu\'{e} celda est\'{a} ocupada y por qu\'{e} entidad (id de enemigo, id de disparo, id de jugador). Es el \'{a}rbitro de colisiones: cuando un disparo se mueve, \texttt{Espacio} consulta la matriz y, si hay colisi\'{o}n, notifica a los observadores. Tambi\'{e}n gestiona la puntuaci\'{o}n y las condiciones de victoria/derrota.
 
    \item \textbf{\texttt{JugadorBueno}} (Singleton + Observer) --- Delega la creaci\'{o}n de la nave a NaveFactory y act\'{u}a como intermediario entre los eventos de teclado y la nave. Observa a \texttt{Espacio} para recibir la se\~{n}al de eliminaci\'{o}n de un proyectil (mensaje tipo 19).
 
    \item \textbf{\texttt{Nave4}} --- Subclase concreta de \texttt{Naves}. Crea una nave distinta de las 3 anteriores, pero el funcionamiento y construccion es idéntico.
 
    \item \textbf{\texttt{Enemigo}} --- Tiene un \texttt{id} \'{u}nico para identificarlo en la matriz y en la vista.
 
    \item \textbf{\texttt{FlotaEnemigos}} (Singleton + Observer) --- Gestiona la lista de enemigos activos. Observa a \texttt{Espacio} para recibir la se\~{n}al de eliminaci\'{o}n de un enemigo (mensaje tipo 18) y la de inicializaci\'{o}n de la flota (mensaje tipo 20).
 
    \item \textbf{\texttt{TimerDisparo}} (Singleton) --- Tick cada 50ms, llama a \url{JugadorBueno.actualizarDisparos()}.
    \item \textbf{\texttt{TimerEnemigo}} (Singleton) --- Tick cada 200ms, llama a \url{FlotaEnemigos.moverEnemigos()}.
\end{itemize}
 
Con esta distribuci\'{o}n, \texttt{Espacio} ya no necesita conocer los detalles de c\'{o}mo se mueven las naves ni c\'{o}mo se crean los disparos. Cada clase tiene una responsabilidad clara y acotada.
\end{cajafinal}
 
% --------------------------------------------------
\subsection{La Capa Vista}
 
\begin{cajasprint1}
Solo exist\'{i}an dos vistas: \texttt{StartFrame} (pantalla de inicio con fondo negro, t\'{i}tulo en verde y mensaje \guillemotleft Pulsa SPACE para jugar\guillemotright) y \texttt{MainFrame} (cuadr\'{i}cula 100x60 con fondo negro). Cada celda era un \texttt{JLabel} con color de fondo: magenta para el jugador, rojo para el enemigo, amarillo para el disparo. El fin de partida se mostraba con un \texttt{JLabel} superpuesto en la parte inferior del propio \texttt{MainFrame}, sin cambiar de pantalla.
\end{cajasprint1}
 
\vspace{0.5em}
\textbf{¿Qu\'{e} motiv\'{o} el cambio?} La pantalla de fin de partida integrada en \texttt{MainFrame} era poco elegante y no permit\'{i}a mostrar la puntuaci\'{o}n ni ofrecer al jugador la opci\'{o}n de volver al men\'{u} de forma limpia. Por ello, se decidi\'{o} a\~{n}adir un frame nuevo al acabar la partida en la versi\'{o}n final. Cabe destacar que esta idea solo se implemento en el Sprint 3, donde se nos ocurrio la idea de la puntuación.
 
\vspace{0.5em}
 
\begin{cajafinal}
Se a\~{n}ade una tercera vista, \texttt{FinalFrame}, y se enriquecen las existentes:
 
\begin{itemize}
    \item \textbf{\texttt{StartFrame}} --- A\~{n}ade selecci\'{o}n de nave (teclas 1/2/3/4), imagen de fondo con \texttt{JLayeredPane} e instrucciones detalladas. La nave elegida se almacena localmente como \texttt{tipoNavePendiente} y solo se confirma al pulsar SPACE, momento en que se llama a JugadorBueno.getJugadorBueno().inicializar(tipoNavePendiente).
 
    \item \textbf{\texttt{MainFrame}} --- El fondo pasa a ser una imagen escalada con \texttt{JLayeredPane}. Los colores de las celdas ya no son \'{u}nicos: cada tipo de nave tiene su propio color (verde para Nave1, azul para Nave2, morado para Nave3, amarillo para Nave4). Se a\~{n}ade un \texttt{JLabel} de puntuaci\'{o}n actualizado por el mensaje tipo 22. Al recibir los mensajes tipo 7 (Game Over) u 8 (Victoria), se desregistra como observer, se oculta y crea un \texttt{FinalFrame}.
 
    \item \textbf{\texttt{FinalFrame}} --- Nueva pantalla de fin con imagen de fondo, mensaje \guillemotleft VICTORIA\guillemotright{} o \guillemotleft GAME OVER\guillemotright{} en color diferente, puntuaci\'{o}n obtenida y opci\'{o}n de volver al men\'{u} pulsando SPACE.
\end{itemize}
 
\textbf{Nota importante:} el hecho de que \texttt{MainFrame} distinga entre m\'{a}s de veinte tipos de mensaje no implica que contenga l\'{o}gica de negocio. La vista simplemente recibe un array de enteros y, seg\'{u}n el tipo, pinta o borra celdas. No toma ninguna decisi\'{o}n sobre el estado del juego.
\end{cajafinal}
 
% --------------------------------------------------
\subsection{La Capa Controlador}
 
En todas las versiones los controladores son clases internas privadas llamadas \texttt{Controller} dentro de cada frame, implementando \texttt{KeyListener}. Esto es correcto desde el punto de vista del patr\'{o}n MVC: el controlador pertenece a la vista pero no es la vista. El flujo com\'{u}n es: el frame crea una instancia del controlador, la registra con \texttt{addKeyListener(controller)} y solicita el foco con \texttt{requestFocusInWindow()}. Solo \texttt{keyPressed} contiene l\'{o}gica real; \texttt{keyReleased} y \texttt{keyTyped} quedan vac\'{i}os.
 
\begin{cajasprint1}
\begin{itemize}
    \item \textbf{Controller de \texttt{StartFrame}}: solo detectaba SPACE y llamaba a
    \url{Espacio.getEspacio().cambiarAMain()}.
    Controller de \texttt{MainFrame}: traduc\'{i}a las cuatro flechas en llamadas a
    \url{Espacio.getEspacio().moverJugador(dx, dy)} y SPACE en
    \texttt{Espacio.getEspacio().disparar()}.
\end{itemize}
El controlador hablaba directamente con \texttt{Espacio}, que era el \'{u}nico punto de entrada al modelo.
\end{cajasprint1}
 
\vspace{0.5em}
\textbf{¿Qu\'{e} motiv\'{o} el cambio?} Al introducir \texttt{JugadorBueno} como singleton que encapsula la nave y sus acciones, \texttt{Espacio} dej\'{o} de ser el punto de entrada para las acciones del jugador. El controlador deb\'{i}a actualizarse para reflejar esta nueva distribuci\'{o}n de responsabilidades. El controlador solo fue cambiado en el segundo Sprint, debido al cambio de rol de la clase Espacio y la creación de la nueva clase JugadorBueno. De tal forma, la versión del segundo Sprint, en cuanto al Controller respecta, es la versión final de nuestro proyecto.
 
\vspace{0.5em}
 
\begin{cajafinal}
\begin{itemize}
    \item \textbf{Controller de \texttt{StartFrame}}: detecta las teclas 1/2/3/4 para actualizar \url{tipoNavePendiente} (solo en la vista, sin tocar el modelo) y SPACE para llamar a \url{JugadorBueno.getJugadorBueno().inicializar(tipoNavePendiente)}. El modelo ya no necesita saber de la transici\'{o}n de pantalla.
    \item \textbf{Controller de \texttt{MainFrame}}: las flechas llaman a \url{JugadorBueno.getJugadorBueno().mover(dx, dy)}, SPACE llama a \url{JugadorBueno.getJugadorBueno().disparar()}, y la tecla M llama a \url{JugadorBueno.getJugadorBueno().cambiarTipoDisparo()}. El controlador ya no llama a \texttt{Espacio} directamente para las acciones de juego.
    \item \textbf{Controller de \texttt{FinalFrame}}: detecta SPACE y llama al m\'{e}todo privado \url{volverAlMenu()}, que desregistra el observer, oculta la ventana y crea un nuevo \texttt{StartFrame}.
\end{itemize}
 
\textbf{Importante:} en ninguna versi\'{o}n el controlador verifica si la partida ha terminado. Esa validaci\'{o}n ocurre en el modelo (\texttt{Espacio}), que es quien tiene autoridad sobre el estado del juego.
\end{cajafinal}
 
% --------------------------------------------------
\subsection{Diagrama UML}
 
% (Insertar aqu\'{i} el diagrama UML cuando est\'{e} disponible)
 
% ==================================================
\subsection{Patr\'{o}n Observer: Notificaci\'{o}n y Actualizaci\'{o}n}
 
\subsection{M\'{e}todo \texttt{notify}}
 
El m\'{e}todo \texttt{notify} pertenece al \textbf{observable} (el modelo). Su funci\'{o}n es avisar a todos los observadores registrados de que el estado ha cambiado. El observable no necesita conocer los detalles internos de sus observadores: simplemente los recorre y llama a su \texttt{update}, manteniendo un bajo acoplamiento entre las clases.
 
En Java, se implementa siempre con la secuencia obligatoria:
\begin{lstlisting}
setChanged();                            // Marca el observable como modificado
notifyObservers(new int[]{tipo, ...});   // Notifica y pasa datos
\end{lstlisting}
 
Sin \texttt{setChanged()}, \texttt{notifyObservers()} no propaga nada.
 
\subsection{M\'{e}todo \texttt{update}}
 
El m\'{e}todo \texttt{update} pertenece a los \textbf{observadores} (las vistas y ciertos componentes del modelo). Cada observador define aqu\'{i} c\'{o}mo reacciona al cambio. Recibe el objeto observable y el argumento pasado por \texttt{notifyObservers}.
 
En este proyecto, el argumento es siempre un \texttt{int[]}: el primer elemento identifica el tipo de notificaci\'{o}n y los siguientes contienen los datos necesarios (coordenadas, ids, puntuaci\'{o}n). 
 
\subsection{Sistema de notificaciones en Space Invaders}
 
\begin{cajasprint1}
El sistema de mensajes era simple: 10 tipos (0--9), todos gestionados directamente en \texttt{Espacio}. Cada tipo correspond\'{i}a a una acci\'{o}n concreta y visible en la cuadr\'{i}cula:
 
\begin{itemize}
    \item \textbf{Tipo 0} --- Movimiento del jugador: \texttt{\{0, oldX, oldY, newX, newY\}}
    \item \textbf{Tipo 1} --- Nuevo disparo: \texttt{\{1, newX, newY\}}
    \item \textbf{Tipo 2} --- Movimiento del disparo: \texttt{\{2, oldX, oldY, newX, newY\}}
    \item \textbf{Tipo 3} --- Disparo sale del tablero: \texttt{\{3, oldX, oldY\}}
    \item \textbf{Tipo 4} --- Movimiento de enemigo: \texttt{\{4, oldX, oldY, newX, newY\}}
    \item \textbf{Tipo 5} --- Colisi\'{o}n disparo-enemigo: \texttt{\{5, disparoX, disparoY, enemigoX, enemigoY\}}
    \item \textbf{Tipo 6} --- Inicializaci\'{o}n del juego: 
    \texttt{\{6, jugadorX, jugadorY, enemigoX, enemigoY\}}
    \item \textbf{Tipo 7} --- Game Over: \texttt{\{7\}}
    \item \textbf{Tipo 8} --- Victoria: \texttt{\{8\}}
    \item \textbf{Tipo 9} --- Cambio a MainFrame: \texttt{\{9\}}
\end{itemize}
 
Cada notificaci\'{o}n correspond\'{i}a a una \'{u}nica celda (o par de celdas), lo que era suficiente para representar un jugador de un p\'{i}xel y un enemigo de un p\'{i}xel.
\end{cajasprint1}
 
\vspace{0.5em}
\textbf{¿Qu\'{e} motiv\'{o} el cambio?} Al pasar a naves multip\'{i}xel, un solo par de coordenadas ya no era suficiente. Adem\'{a}s, los disparos compuestos (flecha, rombo) ocupan varias celdas simult\'{a}neamente. Fue necesario ampliar el protocolo de mensajes y separar la notificaci\'{o}n de la vista de la notificaci\'{o}n interna de gesti\'{o}n.
 
\vspace{0.5em}
 
\begin{cajafinal}
El sistema de mensajes se ampl\'{i}a a m\'{a}s de veinte tipos. Una distinci\'{o}n clave es entre mensajes \textbf{de vista} (tipos 1--3, 10, 12, 14--17, 21, 22), que \texttt{MainFrame} procesa para pintar o borrar celdas, y mensajes \textbf{de modelo} (tipos 18, 19, 20), que \texttt{FlotaEnemigos} y \texttt{JugadorBueno} procesan para actualizar sus listas internas. Los tipos 7, 8 y 9 son recibidos tanto por la vista como por el modelo.
 
\begin{itemize}
    \item \textbf{Tipo 1} --- Nuevo p\'{i}xel de disparo: \texttt{\{1, x, y\}}
    \item \textbf{Tipo 2} --- Movimiento de p\'{i}xel de disparo: \texttt{\{2, oldX, oldY, newX, newY\}}
    \item \textbf{Tipo 3} --- P\'{i}xel de disparo sale del tablero: \texttt{\{3, oldX, oldY\}}
    \item \textbf{Tipo 7} --- Game Over: \texttt{\{7, puntuacion\}}
    \item \textbf{Tipo 8} --- Victoria: \texttt{\{8, puntuacion\}}
    \item \textbf{Tipo 9} --- Cambio a MainFrame: \texttt{\{9\}}
    \item \textbf{Tipo 10} --- Borrar celda de nave jugadora: \texttt{\{10, x, y\}}
    \item \textbf{Tipo 12} --- Borrar p\'{i}xel de enemigo: \texttt{\{12, x, y\}}
    \item \textbf{Tipo 14} --- Pintar p\'{i}xel de enemigo (color gen\'{e}rico): \texttt{\{14, x, y\}}
    \item \textbf{Tipo 15} --- Pintar p\'{i}xel de Nave1 (verde): \texttt{\{15, x, y\}}
    \item \textbf{Tipo 16} --- Pintar p\'{i}xel de Nave2 (azul): \texttt{\{16, x, y\}}
    \item \textbf{Tipo 17} --- Pintar p\'{i}xel de Nave3 (morado): \texttt{\{17, x, y\}}
    \item \textbf{Tipo 18} --- Eliminar enemigo de la flota (solo modelo): \texttt{\{18, enemigoId\}}
    \item \textbf{Tipo 19} --- Eliminar disparo (solo modelo): \texttt{\{19, disparoId\}}
    \item \textbf{Tipo 20} --- Inicializar flota de enemigos: \texttt{\{20, n\_enemigos\}}
    \item \textbf{Tipo 21} --- Pintar p\'{i}xel de Nave4 (amarillo): \texttt{\{21, x, y\}}
    \item \textbf{Tipo 22} --- Puntuaci\'{o}n actualizada: \texttt{\{22, puntuacion\}}
\end{itemize}
 
N\'{o}tese que en la versi\'{o}n final ya no existe el tipo 4 (movimiento de enemigo como celda \'{u}nica) ni el tipo 6 (inicializaci\'{o}n como dos celdas): ahora el movimiento y la inicializaci\'{o}n de enemigos multip\'{i}xel se notifican p\'{i}xel a p\'{i}xel mediante los tipos 12 y 14.
\end{cajafinal}
 
\subsection{Transici\'{o}n de pantallas}
 
\begin{cajasprint1}
El flujo era lineal y sencillo:
\begin{enumerate}
    \item Usuario pulsa SPACE en \texttt{StartFrame}.
    \item Controller llama a \texttt{Espacio.getEspacio().cambiarAMain()}.
    \item \texttt{cambiarAMain()} inicializa jugador y enemigo, arranca el timer, notifica tipo 9 y luego tipo 6.
    \item \texttt{StartFrame.update()} recibe tipo 9: se desregistra, se oculta y crea \texttt{MainFrame}.
    \item \texttt{MainFrame} recibe tipo 6 y pinta las posiciones iniciales.
    \item Al acabar la partida (tipo 7 u 8), \texttt{MainFrame} muestra un \texttt{JLabel} con el resultado. No hay nueva pantalla.
\end{enumerate}
\end{cajasprint1}
 
\vspace{0.5em}
\textbf{¿Qu\'{e} motiv\'{o} el cambio?} Mostrar el resultado como texto superpuesto era una soluci\'{o}n provisional que no permit\'{i}a mostrar la puntuaci\'{o}n final ni ofrecer al usuario la opci\'{o}n de volver al men\'{u} de forma limpia.
 
\vspace{0.5em}
 
\begin{cajafinal}
El flujo completo de transiciones queda as\'{i}:
 
\textbf{StartFrame \textrightarrow{} MainFrame:}
\begin{enumerate}
    \item Usuario elige nave (teclas 1/2/3/4); solo se actualiza \texttt{tipoNavePendiente} en la vista.
    \item Usuario pulsa SPACE: Controller llama a \url{JugadorBueno.getJugadorBueno().inicializar(tipoNavePendiente)}.
    \item \texttt{JugadorBueno.inicializar()} crea la nave v\'{i}a \texttt{NaveFactory}, llama a \url{nave.inicializar()} que propaga hasta \texttt{Espacio.cambiarAMain()}.
    \item \texttt{Espacio.cambiarAMain()} notifica tipo 9.
    \item \texttt{StartFrame.update()} recibe tipo 9: se desregistra, se oculta y crea \texttt{MainFrame}.
    \item \texttt{JugadorBueno} registra la posici\'{o}n inicial en la matriz y la vista (notificaciones tipo 15/16/17/21).
    \item \texttt{Espacio} notifica tipo 20 para que \texttt{FlotaEnemigos} inicialice la flota.
\end{enumerate}
 
\textbf{MainFrame \textrightarrow{} FinalFrame:}
\begin{enumerate}
    \item \texttt{Espacio} detecta game over o victoria y notifica tipo 7 u 8 con la puntuaci\'{o}n.
    \item \texttt{MainFrame.update()} recibe el tipo: se desregistra como observer, se oculta y crea \texttt{new FinalFrame(victoria, puntuacion)}.
\end{enumerate}
 
\textbf{FinalFrame \textrightarrow{} StartFrame:}
\begin{enumerate}
    \item Usuario pulsa SPACE.
    \item Controller llama a \texttt{volverAlMenu()}: desregistra observer, oculta y crea StartFrame.
\end{enumerate}
 
En ninguna de estas transiciones una vista conoce directamente a otra. Todas se comunican a trav\'{e}s de \texttt{Espacio} (v\'{i}a Observer) o de singletons del modelo (\texttt{JugadorBueno}).
\end{cajafinal}
 
\newpage
 
% ==================================================
\part*{Parte II: Patrones de Dise\~{n}o}
\addcontentsline{toc}{section}{Parte II: Patrones de Dise\~{n}o}
 
% ==================================================
\section{Introducci\'{o}n a los Patrones de Dise\~{n}o}
 
En este proyecto se han utilizado tres patrones de dise\~{n}o para que el c\'{o}digo sea m\'{a}s mantenible y extensible: \textbf{Factory}, \textbf{Strategy} y \textbf{Composite}. Cada uno resuelve un problema concreto, y los tres coexisten y se complementan. La mayor\'{i}a de los patrones ya estaban presentes en concepto desde el Sprint 1, pero se consolidaron y ampliaron al a\~{n}adir nuevas funcionalidades en los sprints posteriores.
 
% ==================================================
\section{Factory: La F\'{a}brica de Naves}
 
\subsection{La Idea}
 
Centralizar la creaci\'{o}n de naves en un \'{u}nico lugar. Si los par\'{a}metros de inicializaci\'{o}n cambian (posici\'{o}n, velocidad), solo hay que modificar un fichero. El resto del c\'{o}digo trabaja con la abstracci\'{o}n \texttt{Naves}, sin conocer las clases concretas.
 
\subsection{Evoluci\'{o}n del patr\'{o}n}
 
\begin{cajasprint1}
No exist\'{i}a \texttt{NaveFactory}. La nave se creaba directamente en \texttt{Espacio.inicializar()}:
\begin{lstlisting}
private void inicializar() {
    jugador = new Jugador(50, 55);
    ...
}
\end{lstlisting}
Solo hab\'{i}a un tipo de nave (\texttt{Jugador}), por lo que no hab\'{i}a necesidad de una f\'{a}brica. Sin embargo, esta l\'{o}gica estaba acoplada a \texttt{Espacio}, y a\~{n}adir una segunda nave habr\'{i}a requerido modificar el modelo central.
\end{cajasprint1}
 
\vspace{0.5em}
\textbf{¿Qu\'{e} motiv\'{o} el cambio?} Al introducir cuatro tipos de nave con geometr\'{i}as y arsenales distintos, era imprescindible un punto centralizado de creaci\'{o}n que desacoplara la elecci\'{o}n del tipo (decisi\'{o}n del usuario en \texttt{StartFrame}) de la instanciaci\'{o}n concreta (responsabilidad de la f\'{a}brica).
 
\vspace{0.5em}
 
\begin{cajafinal}
\texttt{NaveFactory} es un Singleton con un \'{u}nico m\'{e}todo p\'{u}blico \texttt{generate(String pNave)}. Todos los par\'{a}metros de inicializaci\'{o}n (posici\'{o}n inicial x=50, y=55, velocidad=1) est\'{a}n definidos aqu\'{i} de forma centralizada.
 
\begin{lstlisting}
public Naves generate(String pNave) {
    int x = 50, y = 55, velocidad = 1;
    switch (pNave) {
        case "Nave1": return new Nave1(x, y, velocidad);
        case "Nave2": return new Nave2(x, y, velocidad);
        case "Nave3": return new Nave3(x, y, velocidad);
        case "Nave4": return new Nave4(x, y, velocidad);
        default:      return null;
    }
}
\end{lstlisting}
 
El cliente (\texttt{JugadorBueno}) solo depende de \texttt{Naves}. Si se a\~{n}ade una \texttt{Nave5}, solo hay que crear la clase y a\~{n}adir un caso en la f\'{a}brica; ning\'{u}n otro fichero cambia. Esto ya lo demostramos al a\~{n}adir la Nave4 en el Sprint~3.
\end{cajafinal}
 
\subsection{Ventajas del patr\'{o}n Factory}
 
\begin{enumerate}
    \item \textbf{Encapsulaci\'{o}n:} La l\'{o}gica de creaci\'{o}n queda centralizada.
    \item \textbf{Bajo acoplamiento:} El cliente depende solo de la interfaz \texttt{Naves}.
    \item \textbf{Extensibilidad:} A\~{n}adir \texttt{Nave5} no afecta al c\'{o}digo existente.
    \item \textbf{Singleton integrado:} Una \'{u}nica instancia de f\'{a}brica durante toda la ejecuci\'{o}n.
    \item \textbf{Control global:} Posici\'{o}n y velocidad iniciales en un \'{u}nico lugar.
\end{enumerate}
 
% ==================================================
\section{Strategy: Disparos Intercambiables}
 
\subsection{La Idea}
 
Cada nave tiene acceso a diferentes tipos de disparo. Unos tienen munici\'{o}n infinita, otros limitada. Unos son un p\'{i}xel, otros forman figuras compuestas. Strategy permite que cada nave tenga su propio arsenal sin duplicar c\'{o}digo y sin condicionales complejos, y adem\'{a}s permite cambiar de tipo de disparo en tiempo de ejecuci\'{o}n.
 
\subsection{Evoluci\'{o}n del patr\'{o}n}
 
\begin{cajasprint1}
No exist\'{i}a el patr\'{o}n Strategy. El disparo era un objeto \texttt{Disparo} simple (un \'{u}nico p\'{i}xel, municion infinita impl\'{i}cita) creado directamente en \texttt{Jugador}. Solo hab\'{i}a un tipo de disparo posible y no era intercambiable en tiempo de ejecuci\'{o}n.
 
\begin{lstlisting}
public void disparar() {
    if (!disparo.isActivo()) {
        disparo = new Disparo(x, y - 1);
        disparo.setActivo(true);
    }
}
\end{lstlisting}
 
A\~{n}adir un segundo tipo de disparo habr\'{i}a requerido \texttt{if-else} dentro de \texttt{Jugador}, acoplando el comportamiento del disparo con la clase de la nave.
\end{cajasprint1}
 
\vspace{0.5em}
\textbf{¿Qu\'{e} motiv\'{o} el cambio?} Al dise\~{n}ar cuatro naves con arsenales diferentes y la posibilidad de cambiar de tipo de disparo en tiempo de ejecuci\'{o}n (tecla M), era necesario encapsular cada comportamiento de disparo en una clase propia e intercambiable.
 
\vspace{0.5em}
 
\begin{cajafinal}
Se define la interfaz \texttt{StrategyDisparo} con el contrato com\'{u}n:
\begin{lstlisting}
public interface StrategyDisparo {
    String getTipo();
    int getMunicion();   // -1 = infinita
    boolean gastar();
    boolean tieneMunicion();
    Component crearDisparo(int x, int y);
}
\end{lstlisting}
 
Las cuatro estrategias concretas son:
 
\begin{description}[leftmargin=2em, style=nextline]
    \item[\textbf{DisparoPixel}] Municion infinita. Crea un \'{u}nico \texttt{Pixel} en la posici\'{o}n de origen. Disponible en todas las naves como disparo base.
    \item[\textbf{DisparoFlecha}] 30 disparos. Crea un \texttt{Composite} con tres \texttt{Pixel} en forma de V invertida: punta en \texttt{(x, y)}, alas en \texttt{(x-1, y+1)} y \texttt{(x+1, y+1)}. Disponible en Nave1 y Nave3.
    \item[\textbf{DisparoRombo}] 20 disparos. Crea un \texttt{Composite} con cinco \texttt{Pixel} en forma de rombo. Disponible en Nave2 y Nave3.
    \item[\textbf{DisparoDoble}] 25 disparos. Crea un \texttt{Composite} con dos \texttt{Pixel} en paralelo: \texttt{(x-1, y)} y \texttt{(x+1, y)}. Disponible en Nave4.
\end{description}
 
La clase \texttt{Disparo} act\'{u}a como \textbf{contexto} del patr\'{o}n: mantiene la lista de estrategias disponibles para la nave actual y el \'{i}ndice de la estrategia activa. Al cambiar de tipo (m\'{e}todo \texttt{cambiarTipoDisparo()}), recorre la lista en orden circular saltando las que no tienen municion.
 
El arsenal de cada nave se configura en su constructor:
\begin{lstlisting}
// Nave3 tiene los tres tipos de disparo
ArrayList<StrategyDisparo> estrategias = new ArrayList<>();
estrategias.add(new DisparoPixel());
estrategias.add(new DisparoFlecha());
estrategias.add(new DisparoRombo());
setGestorDisparos(new Disparo(estrategias));
\end{lstlisting}
\end{cajafinal}
 
\subsection{Ventajas del patr\'{o}n Strategy}
 
\begin{enumerate}
    \item \textbf{Intercambiabilidad en tiempo de ejecuci\'{o}n:} El jugador cambia entre disparos con la tecla M sin reiniciar nada.
    \item \textbf{Encapsulaci\'{o}n:} Cada estrategia gestiona su propia municion y crea su propia forma.
    \item \textbf{Eliminaci\'{o}n de condicionales:} Sin Strategy habr\'{i}a \texttt{if-else} en cada m\'{e}todo de disparo.
    \item \textbf{Extensibilidad:} Un nuevo disparo (p.ej.\ \texttt{DisparoLaser}) solo requiere implementar la interfaz y a\~{n}adirlo al constructor de la nave que lo use.
    \item \textbf{Diferenciaci\'{o}n de naves:} El arsenal es parte de la identidad de cada nave, no un a\~{n}adido externo.
\end{enumerate}
 
% ==================================================
\section{Composite: P\'{i}xeles, Naves y Disparos}
 
\subsection{La Idea}
 
Una nave est\'{a} hecha de p\'{i}xeles. Un disparo compuesto tambi\'{e}n est\'{a} hecho de p\'{i}xeles. Composite permite tratarlos de forma uniforme: mover una nave mueve autom\'{a}ticamente todos sus p\'{i}xeles; comprobar una colisi\'{o}n es simplemente comparar posiciones de p\'{i}xeles.
 
\subsection{Evoluci\'{o}n del patr\'{o}n}
 
\begin{cajasprint1}
No exist\'{i}a el patr\'{o}n Composite. Cada nave y cada disparo era un \'{u}nico p\'{i}xel representado por coordenadas \texttt{x} e \texttt{y}. La detecci\'{o}n de colisiones comparaba directamente esas coordenadas con una ventana de tolerancia de 2 p\'{i}xeles para compensar la discreci\'{o}n del movimiento.
 
Esto funcionaba perfectamente para el dise\~{n}o del Sprint 1, pero no era extensible: representar una nave de 10 p\'{i}xeles con este modelo habr\'{i}a requerido gestionar 10 pares de coordenadas manualmente en cada operaci\'{o}n.
\end{cajasprint1}
 
\vspace{0.5em}
\textbf{¿Qu\'{e} motiv\'{o} el cambio?} Al dise\~{n}ar naves con formas geom\'{e}tricas propias (entre 4 y 10 p\'{i}xeles) y disparos con formas distintas (1, 3 o 5 p\'{i}xeles), era necesario una estructura que permitiera tratar un conjunto de p\'{i}xeles como una \'{u}nica entidad manejable.
 
\vspace{0.5em}
 
\begin{cajafinal}
Se define la interfaz \texttt{Component} con el contrato com\'{u}n para todos los elementos del tablero:
\begin{lstlisting}
public interface Component {
    void mover(int dx, int dy, int tipoNave);
    int getRefX();
    int getRefY();
    boolean isActivo();
    void setActivo(boolean b);
    // + metodos default para notificaciones
}
\end{lstlisting}
 
La jerarqu\'{i}a tiene dos implementaciones:
 
\begin{description}[leftmargin=2em, style=nextline]
    \item[\textbf{Pixel}] La hoja del \'{a}rbol. Representa una \'{u}nica celda del tablero con posici\'{o}n \texttt{(x,y)} y un flag \texttt{esProyectil}. Los p\'{i}xeles de nave solo registran su posici\'{o}n al moverse. Los p\'{i}xeles de proyectil notifican a \texttt{Espacio} en cada movimiento para que la matriz espejo detecte colisiones, y se desactivan si \texttt{y < 0}.
 
    \item[\textbf{Composite}] El nodo interno. Contiene una lista de \texttt{Component} (normalmente \texttt{Pixel}). Tiene dos modos:
    \begin{itemize}
        \item \textbf{Modo nave} (\texttt{esProyectil = false}): antes de mover, valida que ning\'{u}n p\'{i}xel se saldr\'{i}a del tablero. Si alguno se saldr\'{i}a, no mueve nada. Al mover, desplaza todos los p\'{i}xeles y notifica el movimiento en bloque a \texttt{Espacio}.
        \item \textbf{Modo disparo} (\texttt{esProyectil = true}): se mueve libremente; se considera activo mientras al menos un p\'{i}xel hijo est\'{e} activo.
    \end{itemize}
\end{description}
 
Cada nave define su geometr\'{i}a en el m\'{e}todo abstracto \texttt{construir()}:
\begin{lstlisting}
// Nave1: forma de T invertida, 10 pixeles
@Override
public void construir() {
    anadirComponente(new Pixel(x,     y,     false, -1));
    anadirComponente(new Pixel(x + 2, y,     false, -1));
    anadirComponente(new Pixel(x,     y + 1, false, -1));
    anadirComponente(new Pixel(x + 1, y + 1, false, -1));
    anadirComponente(new Pixel(x + 2, y + 1, false, -1));
    anadirComponente(new Pixel(x - 1, y + 2, false, -1));
    anadirComponente(new Pixel(x,     y + 2, false, -1));
    anadirComponente(new Pixel(x + 1, y + 2, false, -1));
    anadirComponente(new Pixel(x + 2, y + 2, false, -1));
    anadirComponente(new Pixel(x + 3, y + 2, false, -1));
}
\end{lstlisting}
 
Del mismo modo, cada estrategia de disparo construye su forma a\~{n}adiendo p\'{i}xeles a un \texttt{Composite}:
\begin{lstlisting}
// DisparoFlecha: forma de V invertida, 3 pixeles
Composite comp = new Composite(true, idDisparo);
comp.addComponent(new Pixel(x,     y,     true, idDisparo)); // punta
comp.addComponent(new Pixel(x - 1, y + 1, true, idDisparo)); // ala izq
comp.addComponent(new Pixel(x + 1, y + 1, true, idDisparo)); // ala der
\end{lstlisting}
 
La detecci\'{o}n de colisiones ya no usa ventana de tolerancia. La \textbf{matriz espejo} de \texttt{Espacio} registra el \texttt{id} del ocupante de cada celda. Cuando un \texttt{Pixel} proyectil se mueve a \texttt{(newX, newY)}, \texttt{Espacio} consulta \texttt{matriz[newX][newY]}: si hay un enemigo ah\'{i}, genera la notificaci\'{o}n de colisi\'{o}n (tipos 12, 18, 19) con precisi\'{o}n absoluta de un p\'{i}xel.
\end{cajafinal}
 
\subsection{Ventajas del patr\'{o}n Composite}
 
\begin{enumerate}
    \item \textbf{Uniformidad:} P\'{i}xeles individuales y grupos de p\'{i}xeles se tratan exactamente igual.
    \item \textbf{Movimiento autom\'{a}tico:} Mover un \texttt{Composite} mueve todos sus hijos recursivamente.
    \item \textbf{Colisiones precisas:} Al trabajar a nivel de p\'{i}xel, no hace falta ventana de tolerancia.
    \item \textbf{Geometr\'{i}as flexibles:} Cada nave y cada disparo define su forma a\~{n}adiendo p\'{i}xeles; no hay l\'{i}mite de forma ni de tama\~{n}o.
    \item \textbf{Extensibilidad:} A\~{n}adir una nueva forma solo requiere definir los p\'{i}xeles en \texttt{construir()} o en \texttt{crearDisparo()}.
\end{enumerate}
 
% ==================================================
\section{C\'{o}mo Todo Funciona Junto}
 
Los tres patrones se complementan y forman un sistema cohesionado:
 
\begin{itemize}
    \item \textbf{Factory} decide \textit{qu\'{e}} nave se crea y con qu\'{e} par\'{a}metros iniciales.
    \item \textbf{Composite} define \textit{c\'{o}mo} se representa esa nave (y sus disparos) en el tablero.
    \item \textbf{Strategy} determina \textit{qu\'{e}} tipos de disparo puede usar esa nave y c\'{o}mo se construye cada uno.
\end{itemize}
 
El flujo completo de una partida ilustra c\'{o}mo interact\'{u}an: el usuario elige nave en \texttt{StartFrame} y pulsa SPACE. El Controller llama a \texttt{JugadorBueno.inicializar("Nave3")}. \texttt{JugadorBueno} llama a \texttt{NaveFactory.generate("Nave3")} (\textbf{Factory}), que devuelve una \texttt{Nave3} con su \texttt{Composite} de 4 p\'{i}xeles (\textbf{Composite}) y sus tres estrategias de disparo (\textbf{Strategy}). Al disparar, la estrategia activa crea un \texttt{Composite} de disparo con los p\'{i}xeles correspondientes; cada p\'{i}xel, al moverse, notifica a \texttt{Espacio}; \texttt{Espacio} consulta la matriz espejo y, si hay colisi\'{o}n, notifica a todos los observers. \texttt{MainFrame} actualiza la vista; \texttt{FlotaEnemigos} elimina al enemigo; \texttt{JugadorBueno} elimina el proyectil.
 
Todo esto ocurre de forma transparente, con cada clase cumpliendo exactamente su responsabilidad.
 
% ==================================================
\section{Conclusiones Finales}
 
La evoluci\'{o}n de este proyecto entre el Sprint 1 y la versi\'{o}n final ilustra de forma pr\'{a}ctica c\'{o}mo los patrones de dise\~{n}o emergen como respuesta a necesidades reales de extensibilidad:
 
\begin{itemize}[leftmargin=2em]
    \item El Sprint 1 era funcional pero monolit\'{i}co: \texttt{Espacio} concentraba demasiadas responsabilidades y el modelo no escalaba a m\'{u}ltiples naves ni disparos compuestos.
    \item El Sprint 2 introdujo \texttt{JugadorBueno} y \texttt{FlotaEnemigos} como clases, y los tres patrones fundamentales del proyecto: Factory (creaci\'{o}n), Strategy (comportamiento) y Composite (estructura).
    \item La versi\'{o}n final distribuye todas las responsabilidades entre clases con roles bien definidos, conectadas mediante el patron Observer (sincronizaci\'{o}n).
    \item \textbf{MVC + Observer} garantiza que el modelo no conoce las vistas y que las vistas no contienen l\'{o}gica de negocio. Además, el Observer permite la comunicación entre clases internas del modelo sin necesidad de llamadas innecesarias.
    \item \textbf{Factory} centraliza la creaci\'{o}n de naves, facilitando la adici\'{o}n de nuevos tipos sin modificar el resto del sistema.
    \item \textbf{Strategy} encapsula los comportamientos de disparo, permitiendo que sean intercambiables en tiempo de ejecuci\'{o}n.
    \item \textbf{Composite} unifica la representaci\'{o}n de naves y disparos a nivel de p\'{i}xel, habilitando colisiones precisas y geometr\'{i}as arbitrarias.
\end{itemize}
 
El resultado es una arquitectura que no solo funciona, sino que est\'{a} preparada para crecer: a\~{n}adir una nueva nave, un nuevo tipo de disparo o una nueva pantalla son cambios localizados que no afectan al resto del sistema.
 
\end{document}